

\section{Metrics in higher dimensions}

As dimensionality $d$ increases, the geometry of the space changes drastically,
and the volume shifts aggressively toward the boundaries.
- the notorious peak of the orange has more volume that its core (see Appendix)


A. Relative l2 error, for:
IC,TC pred vs true, if analytical known: PDE pred vs true
$$
\frac{
\sum_{i=1}^n(\hat y_i - y_i)^2
}{
\sum_{i=1}^n(y_i)^2
}
$$
where $\hat y_i$ are the predictions and $y_i$ is the ground truth value.
use LHS

B. $L^\infty$ err, for:
IC,TC pred vs true, if analytic known: PDE pred vs true
$$
E_\infty = \max_\Omega\|\hat{y}_i - y_i\|
$$
need to look at the section of the domain where the solution is to have the max value
sample the center via gaussian?

select a few important time slides and compute the errs there
- IC
- in between
- TC

we report total $L^\infty$ error, we mean the max
$$
max(
L^\infty_{pde},\,
L^\infty_{bc},\,
L^\infty_{ic}
)
$$
by relative $L^2$, we mean the average
$$
\frac{
L^2{pde}+
L^2{bc}+
L^2{ic}
}{3}
$$
in this construction the 'relative' keeps its meaning.
Since comparing a zero prediciton with a nonzero solution gives 1.

Together, $L^\infty$ tells us about point error, rel. $L^2$ more about domain error.

%For PINNS during training, we also use:
%- MSE - loss - really needed? - gives same info as normal loss computed during calc...
%- just do rell2 w.r.t IC / TC / analytic


%\begin{itemize}
%\item run 5 times with diff random seeds
%\item sample a large set of points and each loggin step, do:
%\item - compute the residual MSE - effective a loss on this set
%\item - compute the rel2 err for IC
%\item early stopping
%\item - when some loss val reached
%\end{itemize}


In pinns - care more about rel l2 since it is closer to the MSE that pinns are trained to satisfy
some tool tend to minimize l2 more than l inf etc.